\documentclass[11pt,a4paper]{article}

\usepackage{ctex}
\usepackage[T1]{fontenc}
\usepackage[top=2.5cm, bottom=2.5cm, left=2.5cm, right=2.5cm]{geometry}
\usepackage{booktabs}
\usepackage{array}
\usepackage{enumitem}
\usepackage{tabularx}
\usepackage{amssymb}
\usepackage{xcolor}
\definecolor{primary}{HTML}{1a1a2e}
\definecolor{accent}{HTML}{3b82f6}
\definecolor{success}{HTML}{22c55e}
\definecolor{danger}{HTML}{ef4444}
\definecolor{warning}{HTML}{f59e0b}
\definecolor{graytext}{HTML}{6b7280}
\usepackage[hidelinks]{hyperref}
\hypersetup{colorlinks=true, linkcolor=primary, urlcolor=accent}
\usepackage{titlesec}
\titleformat{\section}{\Large\bfseries\color{primary}}{}{0em}{}[\titlerule]
\titleformat{\subsection}{\large\bfseries\color{primary}}{}{0em}{}
\usepackage{fancyhdr}
\pagestyle{fancy}
\fancyhf{}
\fancyhead[L]{\small\color{graytext} 企业级 AI Agent 应用商店}
\fancyhead[R]{\small\color{graytext} TRD}
\fancyfoot[C]{\small\color{graytext} \thepage}
\renewcommand{\headrulewidth}{0.4pt}
\usepackage{tikz}
\usetikzlibrary{shapes.geometric, arrows.meta, positioning}

\title{\textbf{TRD：技术需求文档}\\[0.5em]\large Technical Requirements Document}
\author{万能产品小助手（PM AI）\\魏子政 审阅}
\date{2026-05-06 / 版本 v1.1}

\begin{document}

\maketitle

\begin{center}
\begin{tabular}{>{\bfseries}p{3cm} p{10cm}}
\toprule
文档类型 & TRD（Technical Requirements Document） \\
定位 & 技术需求文档 -- 架构怎么选？接口怎么定义？数据库怎么设计？ \\
核心问题 & 技术层面的详细方案 \\
前置文档 & RFC（技术提案） \\
\bottomrule
\end{tabular}
\end{center}

\vspace{1em}

\section{系统架构}

\begin{center}
\begin{tikzpicture}[
  node distance=0.8cm,
  box/.style={rectangle, rounded corners, draw=primary, fill=primary!5, minimum width=2.8cm, minimum height=0.7cm, align=center, font=\small},
  db/.style={cylinder, draw=accent, fill=accent!5, shape border rotate=90, minimum width=1.5cm, minimum height=0.8cm, align=center, font=\small, aspect=0.25},
  arrow/.style={-{Stealth[length=2.5mm]}, thick, graytext}
]
  \node[box] (fe) {Vue 3 + Element Plus\\（SPA，黑白灰风格）};
  \node[box, right=2cm of fe] (be) {FastAPI\\（async，JWT 认证）};
  \node[db, below left=1.2cm and 0.5cm of be] (pg) {PostgreSQL};
  \node[db, below=1.2cm of be] (os) {OpenSearch};
  \node[db, below right=1.2cm and 0.5cm of be] (minio) {MinIO};
  \node[box, right=2cm of be] (scan) {三层扫描\\Semgrep+ClamAV+LLM};

  \draw[arrow] (fe) -- node[above, font=\footnotesize]{HTTP/REST} (be);
  \draw[arrow] (be) -- (pg);
  \draw[arrow] (be) -- (os);
  \draw[arrow] (be) -- (minio);
  \draw[arrow] (be) -- node[above, font=\footnotesize]{上传包} (scan);
  \draw[arrow] (scan) -- node[right, font=\footnotesize]{报告} (pg);
\end{tikzpicture}
\end{center}

\section{API 设计（已实现）}

\subsection{接口规范}

\begin{itemize}[nosep]
  \item RESTful 风格，统一前缀 \texttt{/api/}
  \item 请求/响应格式：JSON
  \item 认证方式：Bearer Token（JWT）
  \item 分页参数：\texttt{page} + \texttt{page\_size}，默认 20 条/页
  \item 列表响应：\texttt{\{items: [...], total, page, page\_size\}}
  \item 错误响应：\texttt{\{detail: "描述", code: "ERROR\_CODE"\}}
\end{itemize}

\subsection{认证接口}

\begin{table}[h]
\centering
\begin{tabularx}{\linewidth}{lllX}
\toprule
\textbf{方法} & \textbf{路径} & \textbf{说明} & \textbf{权限} \\
\midrule
POST & /api/auth/register & 注册（username + password + email） & 公开 \\
POST & /api/auth/login & 登录（username + password） & 公开 \\
POST & /api/auth/send-code & 发送邮箱验证码（返回 6 位码） & 公开 \\
POST & /api/auth/verify-email & 验证邮箱（email + code），验证通过自动登录 & 公开 \\
\bottomrule
\end{tabularx}
\end{table}

\subsection{Skill 接口}

\begin{table}[h]
\centering
\begin{tabularx}{\linewidth}{lllX}
\toprule
\textbf{方法} & \textbf{路径} & \textbf{说明} & \textbf{权限} \\
\midrule
GET & /api/skills & Skill 列表（category + sort + status + page） & 登录 \\
POST & /api/skills/upload & 上传 Skill zip 包 & 登录 \\
GET & /api/skills/mine & 我的 Skill（所有状态） & 登录 \\
GET & /api/skills/\{id\} & Skill 详情（含扫描结果） & 登录 \\
GET & /api/skills/\{id\}/download & 下载 zip（MinIO 预签名 URL） & 登录 \\
POST & /api/skills/\{id\}/install & 安装到 OpenClaw & 登录 \\
POST & /api/skills/\{id\}/offline & 下架（published $\to$ offline） & 管理员 \\
POST & /api/skills/\{id\}/resubmit & 重新提交（rejected $\to$ pending\_review） & 作者 \\
POST & /api/skills/\{id\}/republish & 重新上架（offline $\to$ pending\_review） & 管理员 \\
GET & /api/skills/admin/all & 管理员查看所有 Skill（状态筛选） & 管理员 \\
\bottomrule
\end{tabularx}
\end{table}

\subsection{审批接口}

\begin{table}[h]
\centering
\begin{tabularx}{\linewidth}{lllX}
\toprule
\textbf{方法} & \textbf{路径} & \textbf{说明} & \textbf{权限} \\
\midrule
GET & /api/reviews & 待审批列表（含三层扫描结果） & 管理员 \\
POST & /api/reviews/\{id\}/approve & 通过审批 & 管理员 \\
POST & /api/reviews/\{id\}/reject & 驳回审批（comment） & 管理员 \\
\bottomrule
\end{tabularx}
\end{table}

\section{数据库设计（已实现）}

\begin{table}[h]
\centering
\begin{tabularx}{\linewidth}{llX}
\toprule
\textbf{表名} & \textbf{用途} & \textbf{关键字段} \\
\midrule
users & 用户表 & id, username, email, email\_verified, hashed\_password, role, is\_active, created\_at \\
skills & Skill 表 & id, name, description, version, author\_id, status, category, package\_url, offline\_comment, created\_at \\
skill\_versions & 版本表 & id, skill\_id, version, package\_url, changelog, created\_at \\
scan\_results & 扫描结果表 & id, skill\_id, scan\_type, result (JSONB), passed, created\_at \\
reviews & 审批记录表 & id, skill\_id, reviewer\_id, decision, comment, created\_at \\
\bottomrule
\end{tabularx}
\end{table}

\subsection{Skill 状态流转}

\begin{verbatim}
scanning --> pending_review --> published（上架）
                              \-> rejected（驳回）

published --> offline（管理员下架，需填原因）
offline --> pending_review（管理员重新上架）
rejected --> pending_review（开发者重新提交）
\end{verbatim}

\subsection{其他接口}

\begin{table}[h]
\centering
\begin{tabularx}{\linewidth}{lllX}
\toprule
\textbf{方法} & \textbf{路径} & \textbf{说明} & \textbf{权限} \\
\midrule
GET & /api/scans/\{id\} & 查看单个 Skill 的扫描报告 & 登录 \\
GET & /api/workflow/steps & 获取工作流步骤定义（静态） & 登录 \\
GET & /api/workflow/overview & 获取各状态 Skill 数量统计 & 登录 \\
GET & /health & 健康检查 & 公开 \\
\bottomrule
\end{tabularx}
\end{table}

\section{前端路由（已实现）}

\begin{table}[h]
\centering
\begin{tabularx}{\linewidth}{llX}
\toprule
\textbf{路径} & \textbf{组件} & \textbf{权限} \\
\midrule
/ & 根路由 & 未登录 $\to$ /explore，已登录 $\to$ /dashboard \\
/login & LoginView & 公开 \\
/register & RegisterView & 公开 \\
/verify-email & VerifyEmailView & 公开 \\
/explore & ExploreView & 公开（搜索+分类+排序） \\
/explore/:id & SkillDetailView & 公开（含下载/安装） \\
/submit & SubmitView & 登录 \\
/my-apps & MyAppsView & 登录（所有状态+重新提交） \\
/review & ReviewView & 管理员（审批工作台） \\
/admin/apps & AdminAppsView & 管理员（应用管理+上下架） \\
/dashboard & DashboardView & 登录（KPI 统计） \\
\bottomrule
\end{tabularx}
\end{table}

\section{部署架构}

\begin{itemize}[nosep]
  \item Docker Compose 编排：PostgreSQL + MinIO + OpenSearch + ClamAV
  \item 后端：uvicorn（开发）/ gunicorn（生产）
  \item 前端：Vite dev server（开发）/ nginx 托管 build 产物（生产）
  \item 数据卷持久化：PostgreSQL、MinIO、OpenSearch
  \item 后端启动时自动补齐数据库新字段（seed.py \_ensure\_schema\_columns）
\end{itemize}

\section{工具链约束}

\begin{itemize}[nosep]
  \item \textbf{开发工具链}：OpenClaw（PM AI）+ Cursor（开发 AI）
  \item \textbf{拍板机制}：有歧义时必须由魏子政拍板
  \item \textbf{文档一致性}：TRD 与 MASTER 不一致时，接口/数据库以实际代码为准，重大变更需拍板
\end{itemize}

\section{文档一致性维护规则}

当 TRD 与 MASTER 总需求文档内容不一致时，按以下规则处理：
\begin{enumerate}[nosep]
  \item \textbf{接口定义、数据库设计等局部不一致}：以 MASTER 总需求文档为主，TRD 同步修正
  \item \textbf{涉及重大流程或业务流转路径的问题}：必须由魏子政拍板决定方向，PM 不得自行假设
  \item \textbf{新增接口或修改接口}：需同步更新 PRD 功能清单和 MASTER 需求表
  \item 每次修正后，同步更新 MASTER、PRD 和 TRD 的版本号
\end{enumerate}

\section*{更新记录}

\begin{tabular}{lll}
\toprule
版本 & 日期 & 变更内容 \\
\midrule
v1.1 & 2026-05-06 & 补充scans/workflow/health接口；Skill列表参数细化；审批接口权限明确；新增文档一致性维护规则 \\
v1.0 & 2026-05-06 & 全面重写：API 接口更新为实际路径（/api/skills/*），新增认证接口（send-code/verify-email）、生命周期接口（offline/resubmit/republish）、管理员接口（admin/all），数据库表新增 email/email\_verified/offline\_comment，路由更新为实际路由 \\
v0.2 & 2026-05-06 & 新增 submit/install 接口 \\
v0.1 & 2026-05-06 & 初始版本 \\
\bottomrule
\end{tabular}

\end{document}
